% BHU PYQ -- Manually transcribed & error-corrected
% Paper: BCH-214 : Cost Accounting
% Exam: B.Com. (Hons.) Semester III Examination, 2016-17

\documentclass[11pt,a4paper]{article}
\usepackage[a4paper,top=2.5cm,bottom=2.5cm,left=2.8cm,right=2.8cm,headheight=14pt]{geometry}
\usepackage{amsmath,amssymb}
\usepackage[T1]{fontenc}
\usepackage[utf8]{inputenc}
\usepackage{lmodern,microtype,fancyhdr,enumitem,hyperref,lastpage,parskip}

\hypersetup{
    colorlinks=true,
    urlcolor=black,
    linkcolor=black,
    pdftitle={BCH-214: Cost Accounting -- BHU B.Com. (Hons.) Sem III 2016-17}
}

\pagestyle{fancy}
\fancyhf{}
\renewcommand{\headrulewidth}{0.4pt}
\renewcommand{\footrulewidth}{0.4pt}
\fancyhead[L]{\small\textit{Banaras Hindu University}}
\fancyhead[R]{\small\textit{BCH-214: Cost Accounting}}
\fancyfoot[C]{\small Page \thepage\ of \pageref{LastPage}}

\newlist{parts}{enumerate}{2}
\setlist[parts,1]{label=(\alph*),leftmargin=*,itemsep=5pt,topsep=3pt}
\setlist[parts,2]{label=(\roman*),leftmargin=*,itemsep=3pt,topsep=2pt}
\newcommand{\pts}[1]{\hfill\textit{[#1]}}

\begin{document}

\begin{center}
  {\large\bfseries BANARAS HINDU UNIVERSITY}\\[2pt]
  {\normalsize B.Com. (Hons.) --- Semester III Examination, 2016-17}\\[6pt]
  \rule{0.8\linewidth}{0.5pt}\\[6pt]
  {\large\bfseries Paper No. BCH-214: Cost Accounting}\\[4pt]
  {\normalsize Time: 3 Hours\hfill Maximum Marks: 70}
\end{center}

\rule{\linewidth}{0.4pt}

\smallskip
\noindent\textit{Instructions: Answer all questions. The figures on the right hand side indicate marks.}

\bigskip

\noindent\textbf{Question 1.}
``Cost accounting system is neither unnecessary nor expensive, rather it is a profitable investment.'' Discuss the statement to show the objectives and importance of cost accounting.\pts{14}

\centerline{\textbf{OR}}

Distinguish between the following:\pts{3.5+3.5+3.5+3.5}
\begin{parts}
  \item Controllable costs and Uncontrollable costs
  \item Direct costs and Indirect costs
  \item Relevant costs and Irrelevant costs
  \item Product costs and Period costs
\end{parts}

\medskip\rule{\linewidth}{0.2pt}

\noindent\textbf{Question 2.}
Write notes on the following:\pts{7+7}
\begin{parts}
  \item Time Keeping
  \item Time Booking
\end{parts}

\centerline{\textbf{OR}}

Give the basis of division of standing charges in calculating machine hour rate.\pts{4}

\smallskip
\noindent\textbf{Question 2(b).}
From the following particulars, compute the machine hour rate:\pts{10}
\begin{itemize}[itemsep=3pt,topsep=2pt]
  \item Purchase price of machine: Rs. 1,90,000
  \item Freight and installation charges: Rs. 10,000
  \item Working Life: 5 Years
  \item Repairs and maintenance: 50\% of Depreciation
  \item Annual Power Expenses @ 25 paise per unit: Rs. 6,000
  \item Charges for an 8-hour day:
    \begin{parts}
      \item Power \hfill Rs. 48
      \item Lubricating Oil \hfill Rs. 40
      \item Consumable stores \hfill Rs. 56
      \item Wages \hfill Rs. 160
    \end{parts}
\end{itemize}

\pagebreak

\noindent\textbf{Question 3.}
How would you treat the following items in a cost sheet?\pts{4+4+2+4}
\begin{parts}
  \item Expenses incurred on the removal and re-erection of machines
  \item Training expenses
  \item Packing expenses
  \item Welfare expenses
\end{parts}

\centerline{\textbf{OR}}

A firm manufactured and sold 5,000 machines in the year 2014. Its Trading and Profit and Loss Account for the year is as follows:

\begin{center}
\small
\begin{tabular}{|l|r||l|r|}
\hline
\textbf{Liabilities} & \textbf{Amount (Rs.)} & \textbf{Assets} & \textbf{Amount (Rs.)} \\
\hline
To Cost of Materials & 4,00,000 & By Sales (5,000 units) & 20,00,000 \\
To Direct wages & 6,00,000 & & \\
To Manufacturing Expenses & 2,50,000 & & \\
To Gross Profit c/d & 7,50,000 & & \\
\hline
\textbf{Total} & \textbf{20,00,000} & \textbf{Total} & \textbf{20,00,000} \\
\hline
\hline
To Salaries & 3,00,000 & By Gross Profit b/d & 7,50,000 \\
To Rent and Taxes & 50,000 & & \\
To Selling Expenses & 1,50,000 & & \\
To General Expenses & 1,00,000 & & \\
To Net Profit & 1,50,000 & & \\
\hline
\textbf{Total} & \textbf{7,50,000} & \textbf{Total} & \textbf{7,50,000} \\
\hline
\end{tabular}
\end{center}

For the year 2015, it is estimated that:
\begin{parts}
  \item The output and sales will be of 6,000 machines.
  \item The prices of material will rise by 25\% in comparison to previous years.
  \item Wages will increase by 12.5\%.
  \item Manufacturing expenses will increase in proportion to the combined cost of materials and wages.
  \item Selling expenses per unit will remain unchanged.
  \item Other expenses will remain unaffected by the rise in output.
\end{parts}
Prepare a statement showing the price at which the machines to be manufactured in 2015 should be marketed so as to earn a profit of 10\% on the selling price.\pts{14}

\medskip\rule{\linewidth}{0.2pt}

\noindent\textbf{Question 4.}
Answer the following questions:\pts{5+9}
\begin{parts}
  \item What are the characteristics of the process costing method?
  \item What do you understand by normal wastage, abnormal wastage, and abnormal gain? State how each should be treated in cost accounts.
\end{parts}

\centerline{\textbf{OR}}

A product passes through three processes: X, Y, and Z. The details of expenses incurred on the three processes are as under:

\begin{center}
\small
\begin{tabular}{|l|c|c|c|}
\hline
\textbf{Particulars} & \textbf{Process X} & \textbf{Process Y} & \textbf{Process Z} \\
\hline
Material used & 5,000 Tons & -- & -- \\
Cost per ton (Rs.) & 200 & -- & -- \\
Manufacturing Expenses (Rs.) & 3,62,500 & 1,10,000 & 45,000 \\
Weight lost in processing & 5\% & 10\% & 20\% \\
Scrap sold (per ton) & Rs. 50 & Rs. 50 & Rs. 50 \\
Scrap Quantity & 250 tons & 200 tons & 100 tons \\
Sale price per ton (Rs.) & 350 & 500 & 800 \\
\hline
\end{tabular}
\end{center}

$\frac{2}{3}$ of process X and $\frac{3}{4}$ of process Y are passed on to the next process and the balance are sold.

You are required to prepare Process Accounts showing cost per ton and profit earned at each process.\pts{14}

\pagebreak

\noindent\textbf{Question 5.}
Answer the following questions:\pts{4+10}
\begin{parts}
  \item What are the characteristics of the contract costing method? Explain.
  \item Explain the following:
    \begin{parts}
      \item Work-in-progress valuation
      \item Work certified and uncertified work
      \item Notional profit
      \item Estimated profit
    \end{parts}
\end{parts}

\centerline{\textbf{OR}}

A contractor commenced a building contract on 1st October, 2013. The contract price is Rs. 4,40,000. The following data pertaining to the contract for the year 2014-15 has been compiled from his books:

\begin{itemize}[itemsep=3pt,topsep=2pt]
  \item \textbf{Opening Balances (April 1, 2014):}
    \begin{parts}
      \item Work-in-progress (certified) \hfill Rs. 53,000
      \item Materials at site \hfill Rs. 2,000
    \end{parts}
  \item \textbf{Expenses incurred during 2014-15:}
    \begin{parts}
      \item Materials issued \hfill Rs. 1,12,000
      \item Wages paid \hfill Rs. 1,08,000
      \item Hire of Plant \hfill Rs. 20,000
      \item Other expenses \hfill Rs. 34,000
    \end{parts}
  \item \textbf{Closing Balances (March 31, 2015):}
    \begin{parts}
      \item Materials at site \hfill Rs. 4,000
      \item Work uncertified \hfill Rs. 8,000
      \item Work certified \hfill Rs. 4,05,000
    \end{parts}
\end{itemize}

The cash received represents 80\% of work certified. It has been estimated that further costs to complete the contract will be Rs. 23,000 including the materials at site as on March 31, 2015.

Prepare the contract account and determine the profit on the contract for the year 2014-15 on the estimated profit basis which has to be credited to the profit and loss account.\pts{14}

\vfill
\rule{\linewidth}{0.4pt}
\begin{center}\small
\href{https://bhu-exam.vercel.app/}{Click here to visit our website}\\[4pt]
\href{https://me-aryan.vercel.app/}{Click here to visit our contributor}\\[4pt]
\href{https://quashan-me.vercel.app/}{Click here to visit our contributor}
\end{center}

\end{document}
